% Part: lambda-calculus
% Chapter: church-rosser
% Section: beta-eta-reduction

\documentclass[../../../include/open-logic-section]{subfiles}

\begin{document}

\olfileid[hi]{lam}{cr}{be}

\olsection{$\beta\eta$-अपचयन}

$\beta\eta$-अपचयन ($\bered$) के लिए चर्च--रॉसर गुण लागू होता है।

\begin{lem}\ollabel{lem:one-par}
  यदि $M \beredone M'$, तो $M \beredpar M'$।
\end{lem}

\begin{proof} 
  $M \beredone M'$ के !!{derivation} पर आगमन द्वारा। यदि $M \bredone M'$,
  ~$\eta$-परिवर्तन द्वारा (अर्थात्
  \olref[syn][eta]{defn:beredone}), तो हम \olref[pbe]{thm:refl} उपयोग करते
  हैं। अन्य स्थितियाँ \olref[b]{lem:one-par} जैसी हैं।
\end{proof}


\begin{lem}\ollabel{lem:par-red}
  यदि $M \beredpar M'$, तो $M \bered M'$।
\end{lem}

\begin{proof} $M \beredpar M'$ के !!{derivation} पर आगमन।

  यदि अंतिम नियम \olref[pbe]{defn:beredpar5} है, तो किन्हीं $x$, $N$,
  $N'$ के लिए $M$, ~$\lambd[x][Nx]$ और $M'$, ~$N'$ है, जहाँ $x \notin
  FV(N)$ और $N \beredpar N'$। अतः पहले $\eta$-परिवर्तन द्वारा
  $\lambd[x][Nx]$ को $N$ में अपचयित कर सकते हैं, और फिर $\beredone$ चरणों
  का वह अनुक्रम लगा सकते हैं जो $N \bered N'$ दिखाता है; यह
  आगमन-परिकल्पना से लागू होता है।
\end{proof}


\begin{lem}\ollabel{lem:str}
  $\bered$ वह सबसे छोटा संक्रामक संबंध है जिसमें $\beredpar$ समाहित है।
\end{lem}

\begin{proof}
  \olref[b]{lem:str} की तरह।
\end{proof}

\begin{thm}\ollabel{thm:cr}
  $\bered$ चर्च--रॉसर गुण संतुष्ट करता है।
\end{thm}

\begin{proof}
  \olref[dap]{thm:str}, \olref[pbe]{thm:cr} और \olref{lem:str} से।
\end{proof}
\end{document}
